% !TeX root = main.tex
% Chapitre 2 --- Algorithmes
% Barème visé : "Algorithmes + description globale" (3 pts), cf. METHODO §1.
% Références : Bitsakis (2018) ch. 2 ; Ackermann et al. (2012), StreamKM++, JEA 17.
%
% Ce chapitre décrit les quatre briques algorithmiques empilées par StreamKM++
% (Lloyd → k-means++ → coreset tree → merge-and-reduce), dans l'ordre où
% chacune sert de fondation à la suivante --- même ordre que le §1.1 de la
% thèse et que notre méthodologie de développement (METHODO §0).
%
% Nécessite dans le préambule du document principal :
%   \usepackage{algorithm}
%   \usepackage{algpseudocode}
%   \usepackage{amsmath}

\chapter{Algorithmes}
\label{ch:algorithmes}

StreamKM++ répond à un problème que le $k$-means classique ne peut pas
traiter tel quel : un flux de points trop volumineux pour tenir en mémoire,
qu'on ne peut lire qu'une seule fois. La solution retenue par Ackermann
et al. (2012) et reprise par Bitsakis (2018) consiste à maintenir en
permanence un \emph{coreset} --- un résumé pondéré de taille bornée $m$,
construit incrémentalement, tel que le clustering optimal du coreset
approxime le clustering optimal des données complètes (Définition 2.2,
thèse). Ce chapitre décrit les quatre briques qui rendent cela possible,
et pour chacune, les décisions prises par le groupe lorsque la thèse
laissait une ambiguïté. Le code correspondant vit dans \texttt{streamkm.core}
(§\ref{sec:algo-lloyd}--\ref{sec:algo-kmeanspp}) et \texttt{streamkm.coreset}
(§\ref{sec:algo-coreset-tree}--\ref{sec:algo-merge-reduce}) ; le portage de ces
briques vers Spark fait l'objet du chapitre~4.

\section{Algorithme de Lloyd pondéré}
\label{sec:algo-lloyd}

L'algorithme de Lloyd (thèse, algorithme 2.1) alterne deux étapes jusqu'à
convergence : affecter chaque point au centre le plus proche, puis recentrer
chaque centre sur la moyenne de son cluster.

\begin{algorithm}[H]
\caption{Lloyd pondéré (thèse, algo. 2.1 --- \texttt{KMeans.lloyd})}
\begin{algorithmic}[1]
\Require points pondérés $P = \{(p_i, w_i)\}_{i=1}^n$, centres initiaux $C = \{c_1,\dots,c_k\}$
\Repeat
  \For{chaque point $p_i \in P$}
    \State affecter $p_i$ au centre $c_j$ le plus proche (distance au carré)
  \EndFor
  \For{chaque cluster $j$}
    \State $c_j \gets \dfrac{\sum_{p_i \in \text{cluster } j} w_i \cdot p_i}{\sum_{p_i \in \text{cluster } j} w_i}$
  \EndFor
\Until{$|\Delta\,\mathrm{coût}| < \mathit{tol}$ ou \textit{maxIter} atteint}
\Ensure centres $C$, coût $= \sum_i w_i \cdot \min_j \lVert p_i - c_j\rVert^2$
\end{algorithmic}
\end{algorithm}

\paragraph{Pourquoi « pondéré » n'est pas un détail.}
La thèse énonce l'algorithme sur des points bruts, tous de poids~1. Mais dans
StreamKM++, Lloyd n'est \emph{jamais} appliqué qu'à un coreset --- un ensemble
où le poids d'un point représente la masse d'un sous-ensemble entier de
points originaux (cf. §\ref{sec:algo-coreset-tree}). Un centre recalculé
comme une moyenne \emph{non} pondérée produirait un résultat silencieusement
faux : le code compile, l'algorithme converge, mais vers un optimum qui
n'a plus de rapport avec les données réelles. C'est le piège classique de
StreamKM++, identifié dès la phase de conception du groupe (METHODO, étape
E1) et vérifié par un test dédié (\texttt{KMeansSpec}, « un point de poids
$w$ équivaut à $w$ copies du même point ») : un point de poids 3 doit
produire exactement le même centre que trois copies de poids 1 du même point.

\paragraph{Un invariant vérifié en continu.}
Le coût de Lloyd décroît de façon monotone à chaque itération --- propriété
classique de l'algorithme, mais aussi le signal de correction le plus direct
en cas de bug d'affectation ou de pondération. \texttt{KMeans.lloyd} le
vérifie par une assertion active en mode debug (\texttt{-ea}), plutôt que de
faire confiance silencieusement à la convergence numérique.

\paragraph{Clusters vides.} Si un cluster ne reçoit aucun point à une
itération, son centre est conservé tel quel plutôt que ré-initialisé. Ce
choix, plus simple que l'alternative classique (ré-amorcer sur le point le
plus éloigné), a aussi l'avantage de préserver la monotonie du coût --- un
cluster vide contribue 0, une ré-initialisation aurait pu faire remonter le
coût et casser l'invariant ci-dessus.

\section{Seeding \texorpdfstring{$k$}{k}-means++}
\label{sec:algo-kmeanspp}

Une initialisation aléatoire de Lloyd converge souvent vers un optimum local
médiocre. Le seeding $k$-means++ (thèse, algorithme 2.2) choisit les centres
initiaux avec une probabilité proportionnelle au carré de leur distance au
centre déjà choisi le plus proche (échantillonnage $D^2$), ce qui répartit
les centres sur les amas plutôt que de les laisser s'accumuler.

\begin{algorithm}[H]
\caption{Seeding $D^2$ pondéré (thèse, algo. 2.2 lignes 1--3 --- \texttt{KMeans.seedPlusPlus})}
\begin{algorithmic}[1]
\Require points pondérés $P$, nombre de centres $k$
\State $c_1 \gets$ un point de $P$ tiré \emph{proportionnellement à son poids} \Comment{voir remarque ci-dessous}
\For{$c \gets 2$ \textbf{à} $k$}
  \State $\mathit{closest}(p) \gets \min_{j<c} \lVert p - c_j\rVert^2$ pour chaque $p \in P$
  \State $c_c \gets$ un point $p \in P$ tiré avec probabilité $\propto w_p \cdot \mathit{closest}(p)$
\EndFor
\Ensure $k$ centres $\{c_1,\dots,c_k\}$
\end{algorithmic}
\end{algorithm}

\paragraph{Écart assumé par rapport à la thèse : premier centre non uniforme.}
L'algorithme 2.2 (ligne 1) tire le tout premier centre \emph{uniformément}
parmi les points. Notre implémentation le tire proportionnellement au poids.
Sur des points bruts (tous de poids 1) les deux tirages coïncident
exactement --- l'écart n'a donc \emph{aucun effet} sur l'algorithme 2.2 pris
isolément. Mais dans le pipeline complet de StreamKM++, $k$-means++ n'est
jamais appliqué qu'à un coreset (thèse, §2.3.4, étape 2), où les poids sont
très hétérogènes : un tirage uniforme y biaiserait systématiquement
l'initialisation vers les régions les moins densément représentées par les
données d'origine. Le tirage proportionnel au poids est donc la
généralisation correcte de la ligne 1 au cas pondéré, pas une simplification.
La même remarque s'applique, pour la même raison, au premier représentant du
coreset tree (§\ref{sec:algo-coreset-tree}).

\paragraph{Complexité.} Le calcul incrémental de $\mathit{closest}(p)$ --- mis
à jour à chaque nouveau centre plutôt que recalculé depuis zéro --- donne un
coût en $O(n \cdot k \cdot d)$ et non $O(n \cdot k^2 \cdot d)$.

\paragraph{Combiner les deux : \texttt{KMeans.fit}.} La thèse (§2.3.4,
étape 2) applique cinq fois la combinaison seeding + Lloyd et retient le
clustering de coût minimal --- c'est \texttt{KMeans.fit(points, k, nRestarts=5)}.
Le §5.3.1 de la thèse montre qu'une seule application suffit en pratique ;
\texttt{nRestarts} reste un paramètre expérimental (E3, ch.~5), pas une
constante algorithmique figée.

\section{Construction du coreset : le coreset tree}
\label{sec:algo-coreset-tree}

Un coreset au sens de la Définition 2.2 (thèse) peut être obtenu par un tirage
$D^2$ direct sur l'ensemble des points (\texttt{AdaptiveCoreset}, algo.~2.3),
mais au prix d'une complexité $\Theta(dnm)$ --- chaque nouveau point du coreset
impose de reparcourir tous les points d'origine. Le \emph{coreset tree}
(thèse, §2.3.3, algo.~2.4) ramène cette complexité à $O(dm^2)$ en organisant
la construction comme un clustering hiérarchique divisif : chaque n\oe{}ud de
l'arbre porte un sous-ensemble de points et un représentant, la racine porte
l'ensemble complet, et chaque scission raffine localement le coreset sans
retoucher les autres branches.

\begin{algorithm}[H]
\caption{Coreset tree (thèse, algo. 2.4 --- \texttt{CoresetTree.build})}
\begin{algorithmic}[1]
\Require points pondérés $P$, taille cible $m \le |P|$
\State $q_1 \gets$ un point de $P$ tiré proportionnellement à son poids \Comment{cf. remarque §\ref{sec:algo-kmeanspp}}
\State racine $\gets$ n\oe{}ud unique contenant $P$, représentant $q_1$
\For{$i \gets 2$ \textbf{à} $m$}
  \State $\ell \gets$ descendre de la racine à une feuille, un fils choisi $\propto$ son \textsc{coût} \Comment{§\ref{par:cost-vs-mass}}
  \State $q_i \gets$ un point de la feuille $\ell$ tiré $\propto$ sa distance au carré au représentant de $\ell$ (éch. $D^2$)
  \State scinder $\ell$ en deux nouvelles feuilles (une par représentant $q_\ell, q_i$)
  \State remonter la mise à jour du coût/de la masse jusqu'à la racine
\EndFor
\Ensure $S = \{q_1,\dots,q_m\}$, chaque $q_i$ pondéré par la masse de sa feuille
\end{algorithmic}
\end{algorithm}

\paragraph{L'ambiguïté centrale de la thèse, et comment on l'a tranchée.}
\label{par:cost-vs-mass}
La thèse attache à chaque n\oe{}ud $v$ un attribut \texttt{weight(v)}, défini
juste avant l'algorithme 2.4 comme $\mathrm{cost}(P_v, q_v)$ --- une somme de
distances au carré --- mais réutilise ensuite ce \emph{même mot} « weight »,
au §2.3.4, pour désigner la masse (le nombre de points) que représente
chaque point du coreset final. Ce sont deux quantités de nature complètement
différente portant le même nom. Notre implémentation les sépare
explicitement dans \texttt{CoresetTree.Node} : le champ \texttt{cost}
(la somme de distances au carré, qui pilote la descente dans l'arbre, étape 1
du §2.3.3) et le champ \texttt{mass} (la masse, qui devient le poids du point
de sortie). Les confondre produit un algorithme qui compile, tourne sans
erreur, et construit un coreset dont les poids sont faux --- un bug purement
silencieux.

Une seconde ambiguïté découle directement de la première : le §2.3.3, étape~1
dit en prose que la descente choisit une feuille « avec une probabilité
proportionnelle à $\mathrm{cost}(P_\ell, q_\ell)$ », mais la ligne~5 de
l'algorithme~2.4 dit « selon leurs \emph{weights} ». En comprenant
\texttt{weight(v)} au sens de sa définition (= cost), les deux formulations
concordent --- c'est cette lecture que confirme le papier original d'Ackermann
et al., et c'est donc le \textbf{coût}, pas la masse, qui pilote la descente
dans \texttt{CoresetTree.descend}. Le raisonnement le confirme aussi a
posteriori : une feuille déjà bien approximée par son représentant (coût
faible) n'a pas besoin d'être raffinée davantage ; utiliser la masse à la
place ferait dériver l'algorithme vers un simple échantillonnage
proportionnel à la densité, ce qui lui ferait perdre sa propriété de coreset.

\paragraph{Invariant vérifié.} La somme des poids en sortie du coreset tree
est \emph{exactement} égale à la somme des poids en entrée --- c'est le test
le plus discriminant de toute la couche \texttt{coreset}
(\texttt{CoresetSpec}, « conserve exactement la masse totale ») : un coreset
dont la masse ne se conserve pas est faux, quelle que soit par ailleurs
la plausibilité de ses centres.

\section{Merge-and-reduce}
\label{sec:algo-merge-reduce}

Le coreset tree construit un coreset à partir d'un ensemble de points connu
d'avance. Pour un \emph{flux}, StreamKM++ ajoute une seconde couche : la
technique merge-and-reduce (thèse, §2.3.4, algo.~2.5), qui maintient un
ensemble de « seaux » (\emph{buckets}) $B_0, B_1, \dots$ et applique le
coreset tree pour les fusionner à la demande, sans jamais avoir besoin de
connaître $n$ à l'avance.

\begin{algorithm}[H]
\caption{Insertion dans le flux (thèse, algo. 2.5 --- \texttt{BucketManager.insert})}
\begin{algorithmic}[1]
\Require point $p$, taille de coreset $m$
\State insérer $p$ dans $B_0$
\If{$B_0$ est plein ($m$ points)}
  \State $S \gets$ ensemble vide, $i \gets 1$
  \State déplacer les points de $B_0$ dans $S$, vider $B_0$
  \While{$B_i$ est non vide}
    \State $S \gets \textsc{CoresetTree}(S \cup B_i,\, m)$ \Comment{réduction à $m$ points}
    \State vider $B_i$ ;\; $i \gets i+1$
  \EndWhile
  \State déplacer $S$ dans $B_i$
\EndIf
\end{algorithmic}
\end{algorithm}

\paragraph{Invariants (thèse, §2.3.4, propriétés 1--2), vérifiés par
\texttt{checkInvariants}.} $B_0$ contient entre 0 et $m$ points bruts (poids
1) ; pour $i \ge 1$, $B_i$ est soit vide, soit contient \emph{exactement} $m$
points, et représente alors $2^{i-1} \cdot m$ points du flux d'origine. Le
nombre de buckets alloués croît en $O(\log(n/m))$ : les niveaux sont créés
à la demande, sans jamais avoir besoin de $n$ a priori --- condition
essentielle pour un flux dont la taille finale est inconnue.

\paragraph{Le coreset final.} À tout instant, l'union de tous les buckets
non vides (\texttt{BucketManager.union}) constitue un résumé valide de tout
ce qui a été vu jusqu'ici, borné par $m \cdot (\lceil\log_2(n/m)\rceil + 2)$
points pondérés (thèse, §2.3.4, clarification~2). Interroger cette union
n'interrompt pas la consommation du flux : c'est ce qui permettra, en
extension future (ch.~4, §4.3), de ré-évaluer le clustering périodiquement
sans tout reconstruire.

\paragraph{Fusion de deux structures indépendantes.} La thèse observe
(§2.3.4) que (1)~l'union de deux $\varepsilon$-coresets d'ensembles disjoints
est un $\varepsilon$-coreset de leur union, et que (2)~un $\varepsilon$-coreset
d'un $\delta$-coreset est un $\big((1+\varepsilon)(1+\delta)-1\big)$-coreset ---
l'erreur se compose \emph{multiplicativement}, pas additivement, à chaque
niveau de fusion. \texttt{BucketManager.merge} implémente exactement cette
opération pour deux managers indépendants (par exemple, deux partitions
distinctes d'un même flux), sans modifier ni l'un ni l'autre. Le résultat est
un manager « collapsé » --- un coreset unique de taille $m$, plus un simple
empilement de buckets valides --- dont toute insertion ultérieure est
délibérément refusée : sémantiquement, un manager fusionné est un résultat,
pas un flux qu'on continuerait à alimenter. Cette opération est le point de
départ direct du portage vers Spark (ch.~4) : elle a exactement la signature
d'un \texttt{combOp} de \texttt{treeReduce}, ce qui explique pourquoi la
couche \texttt{streamkm.spark} n'a eu besoin d'aucune modification de
\texttt{streamkm.coreset} pour distribuer merge-and-reduce sur un cluster.
